\section{Counting arbitrary formulas by interpolation}\label{sec:interpolation-return}
The rational counting routine gives the following bound on arbitrary
formulas. We first count one sparse formula and then sum over the
leaves of Lemma~\ref{lem:sparsify}.

\begin{proposition}\label{prop:interpolation-bound}
For each fixed $k\ge2$, the interpolation algorithm counts the
satisfying assignments of a $k$-CNF on $n\ge20$ declared variables
with $m$ clauses in $(n+m)^{O_k(1)}2^{O_k(\rho(n))}$ bit operations
and $(n+m)^{O_k(1)}$ bits of space.
\end{proposition}

\subsection{Counting one sparse formula}\label{subsec:count-sparse}
Suppose $n\ge20$ and $1\le m\le Bn$, where $B$ is bounded by a fixed
power of $\log n$ depending only on $k$. Construct the initial matrix
by Lemma~\ref{lem:sparse-checks}. This uses
$\exp((\log n)^{O_k(1)})$ bit operations and polynomial space.
Both $P$ and $D_0$ are $(\log n)^{O_k(1)}$, so every candidate in
Lemma~\ref{lem:greedy-hashes} is an instance of
Theorem~\ref{thm:overview-sparse-polynomials} with a single fixed
$c=c(k)$. This applies to discarded candidates as well as the final
system: the counter requires only the number, degree and sparsity of
the equations. Expanding each row and combining monomials costs
polynomial time and space for fixed $k$.

At most $3n(2m)^P+1=\exp((\log n)^{O_k(1)})$ calls suffice, including
one final count. Each call takes
$2^{O_k(n(\log\log n)^2/\sqrt{\log n})}$ bit operations. Since any
fixed power of $\log n$ is $o(n(\log\log n)^2/\sqrt{\log n})$, the
total time, including construction and candidate comparisons, is
\begin{equation}\label{eq:sparse-formula-cost}
 (n+m)^{O_k(1)}2^{O_k(n(\log\log n)^2/\sqrt{\log n})}.
\end{equation}
The final count is $\#F$ by the greedy lemma. A formula with no clauses
or containing an empty clause is handled directly. Candidates are
examined sequentially; only the current rows, a candidate, the best
row and its count need be kept in addition to the counter's polynomial
working space.

\subsection{Returning to a general formula}\label{subsec:general-reduction}

We use disjoint sparsification to partition the satisfying
assignments among sparse formulas on the same $n$ declared
variables. The preceding algorithm counts each of these formulas.
The parameter $\sigma$ controls a tradeoff: making it smaller reduces
the number of formulas but increases their occurrence bound $B$.
Our solver permits $B$ to grow as a fixed power of $\log n$, so we
can let $\sigma$ tend to zero while keeping both the number of
formulas and the cost per formula within the desired time bound.

Lemma~\ref{lem:sparsify}, proved in Appendix~\ref{app:sparsify},
supplies the disjoint leaves and their occurrence bound.

\begin{proof}[Proof of Proposition~\ref{prop:interpolation-bound}]
We first normalize the input in $(n+m)^{O_k(1)}$ time by deleting repeated
literals, tautological clauses, duplicate clauses, and clauses
containing another clause. We retain all $n$ declared variables.
An empty clause gives count zero, and a formula with no clauses
gives count $2^n$. We handle $n<20$ by exhaustive search.
For $n\ge20$, we next select the sparsification parameter. Define
\[
 \lambda_n=\lceil\log_2 n\rceil,\qquad
 \nu_n=\lceil\log_2\lambda_n\rceil,\qquad
 s=\left\lfloor\frac{\sqrt{\lambda_n}}{\nu_n^2}\right\rfloor.
\]
We compute $s$ by integer comparisons, taking the largest
nonnegative integer satisfying $s^2\nu_n^4\le\lambda_n$.
If $s<2$, we use exhaustive search;
$\sqrt{\lambda_n}/\nu_n^2$ tends to infinity, so there are only
finitely many such sizes. Otherwise we apply Lemma~\ref{lem:sparsify}
with $\sigma=1/s$. The tree has at most
$2^{n/s}\operatorname{poly}_k(n,s)$ nodes, and its leaves have
occurrence bound
\[
 B=O_k\bigl((s\log s)^{k-1}\bigr)=(\log n)^{O_k(1)}.
\]
Each leaf uses the original $n$ variables and has at most $Bn$
nonempty clauses. We count each leaf by the sparse-formula algorithm
of Section~\ref{subsec:count-sparse}, whose running time is given by
\eqref{eq:sparse-formula-cost}. Leaves with no clauses or containing
an empty clause are handled directly.
The clause count in that application is the number of clauses in
the leaf; $m$ in the final bound remains the input clause count.

For $s\ge2$, the floor in its definition gives
\[
 \frac ns=O\!\left(\frac{n\nu_n^2}{\sqrt{\lambda_n}}\right)
 =O(\rho(n)).
\]
Thus both the number of leaves and the time to count each leaf
have exponent $O_k(\rho(n))$. Their polynomial factors, together
with normalization and the work at the tree nodes, are bounded by
$(n+m)^{O_k(1)}$. The total number of bit operations is therefore
\[
 (n+m)^{O_k(1)}2^{O_k(\rho(n))}.
\]
The finite exhaustive-search cases are absorbed in the constants.
All parameters, initial checks, and additional rows are computed
by the deterministic procedures above; every counting query uses
the common-zero algorithm of Theorem~\ref{thm:overview-sparse-polynomials}.

To bound the space used by the full algorithm, we traverse the
sparsification tree depth first and count one leaf at a time.
The current path has $\poly_k(n,s)$ nodes, and each normalized
formula has a $\poly_k(n)$-bit encoding. Since
$s\le\sqrt{\lceil\log_2 n\rceil}$, storing this path and counting one
leaf use $\poly_k(n,m)$ bits. Disjointness keeps the sum of the
processed leaf counts between zero and $2^n$, so its accumulator
uses at most $n+1$ bits. Thus the total space is $\poly_k(n,m)$ bits.
\end{proof}

